%% ==================================================
%% Section: Registers of several sites and the depth O(log log(N/eps))
%% Sources: arXiv:2307.01696 (Malz, Styliaris, Wei, Cirac), eqs. (5),
%% (10)-(12) and (16), Lemma 1'(i), and paragraph "Tree-RG circuit with
%% measurements".
%% ==================================================

\section{Registers of several sites and the depth \texorpdfstring{$O(\log\log(N/\varepsilon))$}{O(log log(N/epsilon))}}
\label{sec:ldp_register_tree}

In the tree circuit (\ref{eq:ldp_tree_factorization}) a register above the
lowest level carries $\C^{D^2}$, which is a group of sites once $D^2>d$. This
section groups the sites into registers of $s$ consecutive sites, teleports a
register site by site, and applies a layer of unitaries on pairs of distant
registers in a depth depending only on $s$ and on the number of
nearest-neighbour gates of each unitary. With $s$ a length at which the blocked
tensor is injective, the trees of all the blocks prepare the approximating state
with measurements in depth $O(k)$ for block length $s2^{k+1}$, and Lemma~1'(i)
of \cite{Malz2024Preparation} gives the depth $O(\log\log(N/\varepsilon))$ of
the paragraph ``Tree-RG circuit with measurements'' of the source, for chain
lengths divisible by the block length
\cite{gap:mswc24_tree_measurement_scope}.

Throughout, the sites of the ring are $\{0,\dots,N-1\}$ with addition modulo
$N$, the local labels are $\mathbb Z_d$ with $d\ge1$, and $s\ge1$.

\begin{definition}[Register gates]
    \label{def:ldp_register_gate}
    \lean{MPSPreparation.RegisterGate, MPSPreparation.RegisterGate.offset,
        MPSPreparation.RegisterGate.sites,
        MPSPreparation.RegisterGate.localSites,
        MPSPreparation.RegisterGate.span, MPSPreparation.RegisterGate.zone,
        MPSPreparation.RegisterGate.interior, MPSPreparation.RegisterGate.op,
        MPSPreparation.RegisterGate.localOp,
        MPSPreparation.RegisterGate.offset_lt,
        MPSPreparation.RegisterGate.sites_injective,
        MPSPreparation.RegisterGate.localSites_injective,
        MPSPreparation.RegisterGate.localSites_succ,
        MPSPreparation.RegisterGate.sites_mem_span,
        MPSPreparation.RegisterGate.localSites_mem_span,
        MPSPreparation.RegisterGate.zone_subset_span,
        MPSPreparation.RegisterGate.op_mem_unitary,
        MPSPreparation.RegisterGate.op_mem_supportedOperators,
        MPSPreparation.RegisterGate.add_eq_add_iff,
        MPSPreparation.RegisterGate.two_mul_L_lt}
    \leanok
    \uses{def:ldp_measurement_rounds}
    A \emph{register gate} is a unitary $X$ on $2s$ sites together with a site
    $a$ and $L\ge0$ with $2L+2s\le N$. It acts on the two registers
    $a,\dots,a+s-1$ and $a+2L+s,\dots,a+2L+2s-1$, the first register on the
    first $s$ sites of $X$. Its stretch is $a,\dots,a+2L+2s-1$, its interior the
    $2L$ sites $a+s,\dots,a+s+2L-1$ strictly between the registers, and, for
    $0\le t\le s$, its zone $t$ the $2L$ sites $a+t,\dots,a+t+2L-1$.
\end{definition}

\begin{lemma}[Teleporting registers site by site]
    \label{lem:ldp_register_chains}
    \lean{MPSPreparation.IsZeroOn.permMatrix_cfgPerm_mulVec_of_mapsTo,
        MPSPreparation.RegisterGate.there, MPSPreparation.RegisterGate.back,
        MPSPreparation.RegisterGate.allSites_there_subset,
        MPSPreparation.RegisterGate.allSites_back_subset,
        MPSPreparation.RegisterGate.allSites_there_subset_span,
        MPSPreparation.RegisterGate.allSites_back_subset_span,
        MPSPreparation.RegisterGate.pairSites_there_subset,
        MPSPreparation.RegisterGate.pairSites_back_subset,
        MPSPreparation.RegisterGate.sitePerm_there_apply_of_lt,
        MPSPreparation.RegisterGate.sitePerm_there_apply_self,
        MPSPreparation.RegisterGate.sitePerm_back_eq_inv,
        MPSPreparation.RegisterGate.sitePerm_back_apply_of_lt,
        MPSPreparation.RegisterGate.sitePerm_back_apply_self,
        MPSPreparation.RegisterGate.sitePerm_there_notMem_zone,
        MPSPreparation.RegisterGate.sitePerm_back_notMem_zone,
        MPSPreparation.RegisterGate.thereAll,
        MPSPreparation.RegisterGate.backAll,
        MPSPreparation.RegisterGate.zoneAll,
        MPSPreparation.RegisterGate.valid_thereAll,
        MPSPreparation.RegisterGate.valid_backAll,
        MPSPreparation.RegisterGate.sitePerm_thereAll_apply,
        MPSPreparation.RegisterGate.sitePerm_backAll_apply,
        MPSPreparation.RegisterGate.isZeroOn_chainPerm_thereAll,
        MPSPreparation.RegisterGate.isZeroOn_chainPerm_backAll,
        MPSPreparation.RegisterGate.pairSites_thereAll_subset,
        MPSPreparation.RegisterGate.pairSites_backAll_subset,
        MPSPreparation.RegisterGate.sitePerm_backAll_eq_inv,
        MPSPreparation.RegisterGate.therePerm,
        MPSPreparation.RegisterGate.therePerm_apply,
        MPSPreparation.RegisterGate.therePerm_comp_sites,
        MPSPreparation.TeleportHop.valid_nil,
        MPSPreparation.RegisterGate.thereRounds,
        MPSPreparation.RegisterGate.backRounds,
        MPSPreparation.RegisterGate.sum_depth_thereRounds,
        MPSPreparation.RegisterGate.sum_depth_backRounds,
        MPSPreparation.MeasurementRound.IsImplementationOn.isRoundsImplementationOn,
        MPSPreparation.RegisterGate.isRoundsImplementationOn_thereRounds,
        MPSPreparation.RegisterGate.isRoundsImplementationOn_backRounds}
    \leanok
    \uses{def:ldp_register_gate, thm:ldp_teleportation_round,
        lem:ldp_teleportation_chains, lem:ldp_parallel_chains,
        lem:ldp_site_permutations, lem:ldp_round_implementation_compose}
    Let $g_1,\dots,g_n$ be register gates with pairwise disjoint stretches. For
    $0\le t<s$, the round $t$ teleports, for every gate, the content of the site
    $a+t$ to $a+2L+t$ along a chain of $L$ hops, in one round of depth $2$. The
    rounds $t=s-1,\dots,0$, in this order, implement the product $P$ of the
    transpositions of all the chains on the vectors with $\ket0$ at the
    interiors of all the gates, and their outputs have $\ket0$ at the zones $0$.
    The rounds teleporting back, $t=0,\dots,s-1$, in this order, implement
    $P^{-1}$ on the vectors with $\ket0$ at the zones $0$. The rounds of each
    direction have total depth $2s$.
\end{lemma}
\begin{proof}
    \leanok
    The chain of the site $a+t$ meets the sites $a+t,\dots,a+t+2L$, and before
    it the sites $a+t+1,\dots,a+t+2L$ form the zone $t+1$. Starting from $\ket0$
    at the zone $s$, the interior, the round $t$ moves the content of $a+t$ to
    $a+t+2L$ and leaves $\ket0$ at the zone $t$, since the permutation of the
    chain maps the zone $t+1$ together with $a+t$ onto the zone $t$ together
    with $a+t+2L$. The chains of different gates act on disjoint stretches, so
    the lists of all the gates are valid
    (Lemma~\ref{lem:ldp_parallel_chains}), and by
    Theorem~\ref{thm:ldp_teleportation_round} every round implements the
    permutation of its chains. Composition
    (Lemma~\ref{lem:ldp_round_implementation_compose}) gives the claim; the
    backward rounds run the forward ones in reverse.
\end{proof}

\begin{theorem}[A layer of register gates in constant depth with measurements]
    \label{thm:ldp_register_layer}
    \lean{MPSPreparation.mul_list_prod_mul_of_mul_eq_one,
        MPSPreparation.RegisterGate.rounds,
        MPSPreparation.RegisterGate.sum_depth_rounds,
        MPSPreparation.RegisterGate.localProd,
        MPSPreparation.RegisterGate.isZeroOn_localProd_mulVec,
        MPSPreparation.RegisterGate.permMatrix_mul_localProd_mul_permMatrix,
        MPSPreparation.RegisterGate.isZeroOn_zoneAll_iff,
        MPSPreparation.RegisterGate.isRoundsImplementationOn_rounds,
        MPSPreparation.RegisterGate.isCircuitOn_localProd,
        MPSPreparation.RegisterGate.exists_rounds}
    \leanok
    \uses{lem:ldp_register_chains, lem:ldp_round_implementation_compose}
    Let $g_1,\dots,g_n$ be register gates with pairwise disjoint stretches, each
    unitary $X$ a product of at most $K$ unitaries on neighbouring sites. Then
    measurement rounds of total depth $4s+K+2$ implement the product of the
    gates on the vectors with $\ket0$ at the interiors of all the gates: the
    first registers are teleported next to the second ones as in
    Lemma~\ref{lem:ldp_register_chains}, every $X$ is applied to the $2s$
    consecutive sites $a+2L,\dots,a+2L+2s-1$ by a local circuit of depth $K$,
    and the first registers are teleported back. The depth depends neither on
    the distances nor on the number of gates, and no outcome is post-selected.
\end{theorem}
\begin{proof}
    \leanok
    The unitaries on the consecutive sites act inside the stretches, so they
    keep $\ket0$ at the zones $0$, and their product $G$ is a local circuit of
    depth $K$, applied in one round without hops of depth $K+2$. By
    Lemma~\ref{lem:ldp_register_chains} every output is a multiple of
    $P^{-1}GP\ket v$. Conjugation by $P$ maps the consecutive sites of every
    gate back to its two registers, so $P^{-1}GP$ is the product of the gates.
\end{proof}

\begin{theorem}[Trees of register gates on the blocks]
    \label{thm:ldp_register_tree_rounds}
    \lean{MPSPreparation.embedOp_embedOp, MPSPreparation.blockSite_add_natCast,
        MPSPreparation.blockLayerOp_mul, MPSPreparation.blockLayerOp_one,
        MPSPreparation.blockLayerOp_eq_list_prod, MPSPreparation.regSpan,
        MPSPreparation.regOffset, MPSPreparation.mul_sub_two_add_two_mul,
        MPSPreparation.block_pos, MPSPreparation.regWindow,
        MPSPreparation.regSpan_eq, MPSPreparation.two_le_regSpan,
        MPSPreparation.succ_mul_regSpan_le, MPSPreparation.regOffset_lt,
        MPSPreparation.mul_mul_regSpan_add_lt, MPSPreparation.regWindow_val,
        MPSPreparation.regWindow_injective,
        MPSPreparation.mul_mul_add_ne_of_ne,
        MPSPreparation.regWindow_val_ne_of_interior,
        MPSPreparation.regWindow_injective₂, MPSPreparation.regLevelOp,
        MPSPreparation.regTreeOp, MPSPreparation.block_le,
        MPSPreparation.regGate, MPSPreparation.regLevelGates,
        MPSPreparation.two_mul_regGate_L, MPSPreparation.regGate_a_add,
        MPSPreparation.regGate_sites, MPSPreparation.mem_span_regGate,
        MPSPreparation.mem_interior_regGate,
        MPSPreparation.blockSite_ne_of_subBlock_ne,
        MPSPreparation.mem_regLevelGates,
        MPSPreparation.pairwise_regLevelGates,
        MPSPreparation.list_prod_map_op_regLevelGates,
        MPSPreparation.regWindow_val_ne, MPSPreparation.regInteriorFrom,
        MPSPreparation.regInteriorFrom_succ_subset,
        MPSPreparation.interior_subset_regInteriorFrom,
        MPSPreparation.disjoint_range_sites_interior,
        MPSPreparation.isZeroOn_regInteriorFrom_mulVec,
        MPSPreparation.regCentralSites,
        MPSPreparation.regInteriorFrom_zero_subset,
        MPSPreparation.exists_rounds_regTreeOp_aux,
        MPSPreparation.exists_rounds_blockLayerOp_regTreeOp}
    \leanok
    \uses{thm:ldp_register_layer, lem:ldp_round_implementation_compose}
    Let $k\ge0$, $q=s2^{k+1}$, and cut the ring of $N=Mq$ sites into $M\ge1$
    blocks of $q$ sites, each grouped into $2^{k+1}$ registers of $s$
    consecutive sites. For $0\le j\le k$ and $0\le p<2^j$ the sub-block $p$ of
    depth $j$ consists of the registers $p2^{k+1-j},\dots,(p+1)2^{k+1-j}-1$, and
    a unitary $X_{j,p}$ on $2s$ sites acts on its first and its last register.
    Let $T$ be the product of the depths $k,\dots,0$ of a block, the depth $0$
    applied first. If every $X_{j,p}$ is a product of at most $K$ unitaries on
    neighbouring sites, measurement rounds of total depth $(k+1)(4s+K+2)$
    implement $T$ on every block, $T^{\otimes M}$, on the vectors with $\ket0$
    at every site of every block other than its first and its last $s$ sites.
\end{theorem}
\begin{proof}
    \leanok
    The unitaries of depth $j$ of all the blocks are register gates whose
    stretches are their sub-blocks, which are pairwise disjoint, so
    Theorem~\ref{thm:ldp_register_layer} applies the depth $j$ in depth
    $4s+K+2$ on the vectors with $\ket0$ at the sites strictly inside the
    sub-blocks of depth $j$ between their first and last registers. A unitary
    of depth $j'<j$ acts on the first and the last register of a sub-block of
    depth $j'$, which are the first or the last register of a sub-block of
    depth $j$, so on no site where depth $j$ needs $\ket0$. Induction on the
    depth and Lemma~\ref{lem:ldp_round_implementation_compose} give the claim.
\end{proof}

\begin{lemma}[Amplitudes of trees and of layers of placed unitaries]
    \label{lem:ldp_tree_amplitudes}
    \lean{MPSTensor.unpair, MPSTensor.two_mul_add_lt_two_pow,
        MPSTensor.unpair_apply, MPSTensor.unpair_injective,
        MPSTensor.unpair_surjective, MPSTensor.unpair_bijective,
        MPSTensor.decodeBlock_blockIndexOfList, MPSTensor.decodeBlock_finCongr,
        MPSTensor.decodeBlock_pairRegroupEquiv, MPSTensor.treeAmp,
        MPSTensor.binaryTreeMatrix_apply_unpair, MPSPreparation.placeCfg,
        MPSPreparation.placeCfg_apply, MPSPreparation.placeCfg_comp,
        MPSPreparation.placeCfg_apply_of_notMem, MPSPreparation.eq_placeCfg,
        MPSPreparation.placeCfg_injective, MPSPreparation.injective_window,
        MPSPreparation.disjoint_range_window,
        MPSPreparation.list_prod_embedOp_placeCfg,
        MPSPreparation.mul_apply_extend}
    \leanok
    \uses{def:ldp_binary_mera}
    \begin{enumerate}
        \item The tree map (\ref{eq:ldp_binary_mera_tree}) with finest layer
              $W$ and coarser layers $V_0,V_1,\dots$, at $2^{j+1}$ leaves grouped
              into neighbouring pairs $e_0,\dots,e_{2^j-1}$, is
              $\sum_{e'}\prod_pW(e_p,e'_p)$ times the tree map of the coarser
              layers at the outputs $e'$ grouped into pairs in their turn, the
              sum running over these outputs.
        \item Let $Y_1,\dots,Y_n$ be unitaries on pairwise disjoint windows of
              $m$ sites of a chain, with $Y_p\ket{\iota_p(c)}=\sum_e
                  V_p(e,c)\ket{o_p(e)}$ for injective $\iota_p,o_p$. Then the
              product of the placed unitaries maps the configuration with
              $\iota_p(c_p)$ on the window $p$ and $0$ elsewhere to
              $\sum_e\prod_pV_p(e_p,c_p)$ times the configuration with $o_p(e_p)$
              on the window $p$ and $0$ elsewhere.
    \end{enumerate}
\end{lemma}
\begin{proof}
    \leanok
    The first item is the recursion defining the tree map, with the sum of a
    matrix product written out. For the second, a placed unitary acts on its
    window only and the windows are disjoint, so the product factors over the
    windows; induction on $n$.
\end{proof}

\begin{theorem}[The tree of registers implements the isometry of a block]
    \label{thm:ldp_register_tree_isometry}
    \lean{Matrix.polarIso_submatrix_equiv,
        MPSPreparation.exists_unitary_apply_eq_extend,
        MPSPreparation.extUnitary, MPSPreparation.extUnitary_mem_unitary,
        MPSPreparation.extUnitary_apply, MPSPreparation.append_apply_of_lt,
        MPSPreparation.append_apply_of_le, MPSPreparation.append_injective,
        MPSPreparation.regSpan_eq_two_mul, MPSPreparation.regWindow_two_mul,
        MPSPreparation.regWindow_two_mul_add_one,
        MPSPreparation.regWindow_two_mul_notMem,
        MPSPreparation.regWindow_two_mul_add_one_notMem,
        MPSPreparation.exists_regWindow_eq,
        MPSPreparation.polarIsoMatrix_blockTensor_blockTensor,
        MPSPreparation.iteratedBlockIndex_decodeBlockEquiv_symm,
        MPSPreparation.regInput, MPSPreparation.coarseOutput,
        MPSPreparation.fineOutput, MPSPreparation.regInput_injective,
        MPSPreparation.decodeBlock_two_injective,
        MPSPreparation.coarseOutput_injective,
        MPSPreparation.fineOutput_bijective, MPSPreparation.regLayer,
        MPSPreparation.regTreeGate, MPSPreparation.regTreeGate_mem_unitary,
        MPSPreparation.regTreeGate_apply_of_lt,
        MPSPreparation.regTreeGate_apply_of_le,
        MPSPreparation.regRegisterCfg, MPSPreparation.regLeafCfg,
        MPSPreparation.regRegisterCfg_injective,
        MPSPreparation.regLeafCfg_bijective, MPSPreparation.regLeafCfg_apply,
        MPSPreparation.regRegisterCfg_eq_placeCfg,
        MPSPreparation.placeCfg_regInput_zero,
        MPSPreparation.prod_fin_two_pow_zero,
        MPSPreparation.regTreeOp_apply_regRegisterCfg,
        MPSPreparation.regTreeOp_apply_regLeafCfg,
        MPSPreparation.regTreeOp_apply_blockInputCfg}
    \leanok
    \uses{thm:ldp_tree_factorization, lem:ldp_tree_layers_isometries,
        lem:ldp_unitary_implementing_isometry, lem:ldp_tree_amplitudes,
        thm:ldp_register_tree_rounds, thm:ldpolar_tensor_decomposition}
    Let $A$ be a tensor whose blocked tensor $\mathcal B$ over $s$ sites has an
    injective two-site blocked tensor, let $k\ge0$ and $q=s2^{k+1}$, and let
    $\mathcal B_q=VP$ be the polar decomposition of $A$ blocked over $q$ sites.
    Let a register of $s$ sites carry the bond legs $\C^D$ through an injective
    map of $\{0,\dots,D-1\}$ into its configurations, and $\C^{D^2}$ through an
    injective map of $\{0,\dots,D^2-1\}$ into them. Every layer of the tree
    circuit (\ref{eq:ldp_tree_factorization}) of $\mathcal B$ is an isometry
    from one register to two, and extends to a unitary on the two registers of
    each sub-block of Theorem~\ref{thm:ldp_register_tree_rounds}. The tree $T$
    of these unitaries satisfies
    \begin{align*}
        \bra\tau T\ket{l,0,\dots,0,r}=\bra\tau V\ket{l,r}
    \end{align*}
    for every configuration $\tau$ of the block and all bond indices $l,r$,
    where $\ket{l,0,\dots,0,r}$ carries $l$ on the first and $r$ on the last
    register.
\end{theorem}
\begin{proof}
    \leanok
    By Lemma~\ref{lem:ldp_tree_layers_isometries} the layers of the tree circuit
    of $\mathcal B$ are isometries, and they extend to unitaries
    (Lemma~\ref{lem:ldp_unitary_implementing_isometry}) from their input
    register, the other register in $\ket0$, to their two output registers.
    The unitary of the root reads $l$ and $r$. After the depths $0,\dots,j$ the
    registers of the sub-blocks of depth $j+1$ carry the outputs of these
    layers, with the amplitude of the tree of the coarser layers, by the
    second item of Lemma~\ref{lem:ldp_tree_amplitudes}: the registers written by
    a unitary of depth $j$ are the first and the last register of the two
    halves of its sub-block, on which the two unitaries below it act. At the
    finest depth the block carries the leaves, with the amplitude of the whole
    tree map, which by the first item and
    Theorem~\ref{thm:ldp_tree_factorization} is the isometry of the polar
    decomposition of $\mathcal B$ blocked over $2^{k+1}$ sites, that is, of
    $A$ blocked over $q$ sites.
\end{proof}

\begin{theorem}[The approximating state with measurements in depth $O(k)$]
    \label{thm:ldp_register_tree_approximating_state}
    \lean{MPSPreparation.isZeroOn_regCentralSites_pairLayerOp_mulVec,
        MPSPreparation.exists_isPreparedWithMeasurementRoundsInDepth_approximatingMPVState}
    \leanok
    \uses{thm:ldp_register_tree_isometry, thm:ldp_register_tree_rounds,
        thm:ldp_approximating_state_depth, def:ldp_approximating_state,
        thm:ldp_bounded_unitary_two_site_gates, thm:ldp_chain_depth,
        def:ldp_measurement_rounds}
    For every $s\ge1$ there is $C$, depending only on $d$ and $s$, with the
    following property. Let $A$ be a tensor whose blocked tensor over $s$ sites
    is injective, and let $\sigma\ge0$ with $\Tr\sigma=1$. For every $k\ge0$
    and $M\ge1$, the approximating state of
    Definition~\ref{def:ldp_approximating_state} with blocks of $q=s2^{k+1}$
    sites, on $N=Mq$ sites, is prepared with measurement rounds in depth at most
    $C(k+1)$. No outcome is post-selected.
\end{theorem}
\begin{proof}
    \leanok
    Injectivity gives $D^2\le d^s$, so a register of $s$ sites carries
    $\C^{D^2}$ and $\C^D$, and the two-site blocked tensor of the injective
    blocked tensor is injective. By
    Theorem~\ref{thm:ldp_bounded_unitary_two_site_gates} there is $K$,
    depending on $d$ and $s$, such that every unitary on $2s$ sites is a
    product of at most $K$ unitaries on neighbouring sites. A local circuit of
    depth $K$, applied in one round of depth $K+2$, prepares the pairs
    $\ket\omega$ of the fixed point on the last register of every block and the
    first of the next, from the all-$\ket0$ product state; the other sites keep
    $\ket0$. Theorem~\ref{thm:ldp_register_tree_rounds}, with the unitaries of
    Theorem~\ref{thm:ldp_register_tree_isometry}, then applies the trees of all
    the blocks in depth $(k+1)(4s+K+2)$. The output is
    $T^{\otimes M}\bigotimes\ket\omega=V^{\otimes M}\bigotimes\ket\omega$, the
    approximating state, as in Theorem~\ref{thm:ldp_approximating_state_depth};
    the total depth is at most $(2K+4s+4)(k+1)$.
\end{proof}

\begin{theorem}[Preparation with measurements in depth $O(\log\log(N/\varepsilon))$]
    \label{thm:ldp_loglog_depth}
    \lean{MPSPreparation.exists_isPreparedWithMeasurementRoundsInDepth_approximationError_le,
        MPSPreparation.exists_isPreparedWithMeasurementRoundsInDepth_le_logb_log}
    \leanok
    \uses{thm:ldp_register_tree_approximating_state,
        lem:ldp_approximation_error_slope_threshold,
        thm:normalized_primitive_gauge, lem:ldp_normal_gap,
        lem:ldp_approximating_state_unit_norm,
        lem:ldp_normalized_state_rescaling}
    Let $A$ be normal with $D\ge1$. There are $s\ge1$, $C$, $a>0$ and $b\ge1$,
    depending only on $A$, with the following property. Let $\varepsilon>0$,
    $N\ge1$ and $k\ge0$ be such that $s2^{k+1}$ divides $N$ and
    $s2^{k+1}\ge a\log(N/\varepsilon)+b$. Then a unit vector $\ket\psi$ with
    $\varepsilon(\psi,\phi_N)\le\varepsilon$ is prepared with measurement
    rounds in depth at most $C(k+1)$. If moreover
    $s2^{k+1}\le2(a\log(N/\varepsilon)+b)$, then
    $k\le\log_2(a\log(N/\varepsilon)+b)$, so the depth is at most
    $C(\log_2(a\log(N/\varepsilon)+b)+1)=O(\log\log(N/\varepsilon))$. This is
    the depth of the paragraph ``Tree-RG circuit with measurements''
    of~\cite{Malz2024Preparation} for normal tensors, restricted to the chain
    lengths divisible by $s2^{k+1}$ for some $k$ in this window, for instance
    $N=Ms2^K$ with $K$ large \cite{gap:mswc24_tree_measurement_scope}.
\end{theorem}
\begin{proof}
    \leanok
    Bring $A$ into the gauge~(\ref{eq:ldp_normal_gauge}) as in the proof of
    Theorem~\ref{thm:ldp_depth_upper_bound_divisible}, which multiplies
    $\ket{\phi_N}$ by a nonzero scalar, let $\xi$ be the correlation length of a
    bound $t<1$ on the moduli of the subleading eigenvalues, and let $s$ be a
    length at which the blocked tensor in this gauge is injective. Take $a=\xi$
    and $b=b_0+1$ with $b_0$ from
    Lemma~\ref{lem:ldp_approximation_error_slope_threshold}. With $N=Mq$ and
    $q=s2^{k+1}$, the approximating state with blocks of $q$ sites is a unit
    vector with error at most $\varepsilon$, prepared in depth at most
    $C(k+1)$ by Theorem~\ref{thm:ldp_register_tree_approximating_state}. If
    $s2^{k+1}\le2(a\log(N/\varepsilon)+b)$ then $2^k\le
        a\log(N/\varepsilon)+b$ since $s\ge1$.
\end{proof}
